\section{The Euler layer and ordinary reduction}
\label{sec:layer}

We identify the cofiber of the Euler parameter in the mapping algebra
of Section~\ref{sec:power}. The comparison will be multiplicative and
will respect both ordinary scalar classes and the particular null
used to form the cofiber. These compatibilities allow a later algebra
attachment to determine an actual class in a Koszul object.

Define the additive map $F:I\to\Omega$ by
\[
 F(w_h^{(\lambda)})=w_{h+1}-w_1,
 \qquad F(w_h^{(\delta)})=w_h,
 \qquad F(\tau)=-w_1.
\]
Recall that $q=d_1=2(p-1)$. The ordinary degrees satisfy
\begin{equation}
 \label{layer:degree}
 d(Fi)=p\,d(\lambda)+d(\delta)-qm
 \qquad (i=(\lambda,\delta,m)).
\end{equation}
Here $\lambda$ records the source of the permutation power,
$\delta$ the ordinary scalar shift, and $m$ the Euler shift. The
map $F$ records the resulting ordinary source sphere. In particular,
for $0\leq j<n$,
\[
 F(w_j,0,-1)=w_{j+1},\qquad
 d\bigl(F(w_j,0,-1)\bigr)=p d_j+q=d_{j+1}.
\]
This is the sphere degree required for the next attaching class.
At this index a graded target is still evaluated at $p w_j$, as in
$B(X)$; $F$ specifies the source sphere.

For a graded sequence $X$, set
\[
 \begin{split}
 Q(X)_i&=\map_{\cC}
 \left(S(Fi),\bigl(X_p(p\lambda+\delta)\bigr)_p\right),\\
 \Phi(X)_i&=\map_{\cE}
 \left(W_i,J^\vee\otimes
 \Tr\bigl(X_p(p\lambda+\delta)\bigr)_p\right).
 \end{split}
\]
The stable dual $J^\vee$ is a commutative algebra by the diagonal of
$S(V)$. Thus these are graded algebras when $X$ is. We will compare
$B(X)//t$ with $\Phi(X)$, and then compute $\Phi(X)$ as the reduction
of $Q(X)$ by $\varpi$.

\subsection{The sphere with affine action}

\begin{lemma}[Acyclicity of the Borel sphere]
\label{layer:acyclic-sphere}
For every odd prime $p$, the space
$Z_p^{\mathrm{orb}}=S(V)_{hG_p}$ has the $p$-local homology of a
point. Its stabilized augmentation to the $p$-local sphere is an
equivalence.
\end{lemma}

\begin{proof}
The two-periodic resolution of $C_p$ gives
\[
 H^*(BC_p;\F_p)=\Lambda(\eta)\otimes\F_p[\xi],
 \qquad |\eta|=1,\quad |\xi|=2.
\]
The class $\xi$ is the Bockstein of $\eta$, or equivalently the
reduction of a first Chern class. Under an automorphism of $C_p$,
both classes transform with multiplier weight one, up to replacing
the character by its inverse throughout. As a $C_p$-representation,
$V$ is a sum of $(p-1)/2$ nontrivial complex lines. Its Euler class
is a nonzero scalar times $\xi^{(p-1)/2}$. The sphere-bundle sequence
therefore gives
\[
 H^*(S(V)_{hC_p};\F_p)
 =\Lambda(\eta)\otimes
 \F_p[\xi]/\bigl(\xi^{(p-1)/2}\bigr).
\]
Every positive-degree monomial has multiplier weight strictly between
zero and $p-1$. Passage to $G$ takes invariants under
$\F_p^\times$, whose order is invertible in $\F_p$, so only the
degree-zero class remains.

The $p$-local homology is finitely generated in every degree, since
the group is finite and the sphere is finite dimensional. The
universal coefficient sequence and Nakayama's lemma now imply that
its reduced $\Z_{(p)}$-homology vanishes. The augmentation cofiber is
bounded below; the stable Hurewicz theorem applied to its first
possible nonzero homotopy group makes it contractible.
\end{proof}

\begin{lemma}[Even spherical trivializations]
\label{layer:even-trivializations}
For every even integer $r$, including negative integers, the local
$p$-local sphere of $rV$ over $Z_p^{\mathrm{orb}}$ is trivializable
with fiber degree $r(p-1)$. Consequently there are coherent
$J^\vee$-module equivalences
\begin{equation}
 \label{layer:oriented-W}
 J^\vee\otimes W_i\simeq
 J^\vee\otimes\Tr S(Fi)
 \qquad(i\in I).
\end{equation}
They may be chosen to agree with the given trivial identifications
on the $\delta$ summand, and they remain equivalences after
localization.
\end{lemma}

\begin{proof}
For even $r$, the determinant character of $rV$ is trivial. After
fixing the fiber degree, its classifying map to the classifying space
of sphere automorphisms lifts to the simply connected cover. All
higher homotopy groups of this cover are $p$-local abelian groups.
By Lemma~\ref{layer:acyclic-sphere}, reduced cohomology of
$Z_p^{\mathrm{orb}}$ with every such constant coefficient group
vanishes. Induction along the Eilenberg--Mac Lane fibers of the
Postnikov tower shows that the corresponding based mapping spaces
are contractible. Mapping into the limit gives a based null of the
classifying map. Dualization handles virtual negative multiples as
well.

Maps between free twisted rank-one $J^\vee$-modules are sections of
the corresponding sphere mapping local system over
$S(V)_{hG}$. Indeed, this follows from free-module adjunction and
duality for $S(V)_+$; composition of these maps is pointwise
composition of sections. The trivializations just constructed
therefore give the required module equivalences primewise on
generators. In $W_i$, the twisting exponent of $V$ is
$d(\lambda)-2m$, which is even, and its total underlying dimension
is Equation~\eqref{layer:degree}. Apply
Lemma~\ref{found:picard-grading} in the rational category of
$J^\vee$-modules in $\Ind(\cG)$ to extend the generator
equivalences coherently over $I$. All the required generator objects
are squares. The relative assertion in that lemma keeps the
$\delta$ identifications fixed. Finally apply the symmetric
monoidal localization.
\end{proof}

In particular, free-module adjunction allows the coherent replacement
of $W_i$ by $\Tr S(Fi)$ in the definition of $\Phi(X)$. It remains
to compute maps into $J^\vee$ with trivial coefficients. We retain a
bounded-join orbit term in this computation; replacing it by full
homotopy orbits before taking the ultraproduct would lose the required
uniform information.

\subsection{The partial norm calculation}

For $U=(U_p)_p\in\cG$, define spectra
\[
 \begin{split}
 \operatorname{Part}(U)
 &=\colim_b\prod_{\U}(E_{b+}\otimes U_p)_{hG_p},\\
 \operatorname{Full}(U)
 &=\prod_{\U}U_p^{hG_p}.
 \end{split}
\]
Ultraproducts of spectra in these formulas mean filtered colimits of
products. The first construction uses the fixed finite join stages;
the second takes the full homotopy fixed points at each prime.

\begin{lemma}[The partial and full norm model]
\label{layer:norm-model}
There is a natural cofiber sequence
\begin{equation}
 \label{eq:norm-triangle}
 \operatorname{Part}(U)\longrightarrow
 \operatorname{Full}(U)\longrightarrow
 \map_{\cE}(\one,U).
\end{equation}
The first map is projection from a join stage followed by the norm.
The norm is an equivalence on induced spectra and on their finite
stable extensions. On underlying fibers, fiber inclusion followed
by norm and evaluation is the sum of the group actions.
\end{lemma}

\begin{proof}
For a finite group, the norm is the natural map from homotopy orbits
to homotopy fixed points; see
\cite[Definition~I.1.10 and Example~I.1.11]{NS2018}. One concrete
description uses the group ring spectrum $\mathcal H=S_{(p)}[G]$.
To specify its derived construction, regard $\mathcal H$ as a
$G\times G$-spectrum with action $(g,h)\cdot k=gkh^{-1}$.
It is induced from the trivial sphere along the diagonal subgroup
$\Delta G$. Finite induction agrees with coinduction, so the unit
of restriction--coinduction gives a map from the trivial
$G\times G$-sphere to $\mathcal H$. On underlying spectra this
is the diagonal into the finite product of spheres. Equivalently,
it is the $\mathcal H$-bimodule norm-element map that sums all
group elements.

Take its derived tensor over the right $\mathcal H$-action with
$U_p$. The remaining left action makes the resulting map
from the trivially acted-on spectrum $S_{(p)}\otimes_{\mathcal H}U_p$
to $U_p$ equivariant. Its adjoint is
\[
 S_{(p)}\otimes_{\mathcal H}U_p
 \longrightarrow\map_{\mathcal H}(S_{(p)},U_p).
\]
In particular, composing the inclusion of an underlying fiber with
this map and then evaluating at a point sums the group actions.
On an induced object $\mathcal H\otimes Z$, the orbit and
fixed-point adjunctions identify both sides with $Z$; the diagonal
entry at the identity element identifies the norm with
$\id_Z$. Both functors are exact, so the same
equivalence holds on finite stable extensions of induced objects;
compare \cite[Lemma~I.3.8]{NS2018}.

Apply this to $E_{b+}\otimes U_p$, whose equivariant cellular
filtration consists of free cells. The join triangle, followed by
homotopy fixed points, thus identifies its first term with homotopy
orbits. Take the ultraproduct and then the colimit in $b$.
Equation~\eqref{power:join-formula} identifies the last term, giving
Equation~\eqref{eq:norm-triangle}.
\end{proof}

\begin{lemma}[Partial Thom spectra with arbitrary coefficients]
\label{layer:partial-thom}
Let $Z=(Z_p)_p$ be any sequence of $p$-local spectra, without
connectivity or boundedness assumptions. Let $r=(r_p)_p$ be a
sequence of integers of fixed parity. If $r$ is odd, then
\[
 \operatorname{Part}(T^r\Tr Z)=0.
\]
If $r$ is even, inclusion of a fiber point at the identity vertex of
$E_b$ induces an equivalence
\[
 \prod_{\U}(T_p^{r_p}\otimes Z_p)\xrightarrow{\simeq}
 \operatorname{Part}(T^r\Tr Z),
\]
where the source has forgotten the group action.
\end{lemma}

\begin{proof}
First omit $Z_p$. The partial Thom spectrum
$(E_{b+}\otimes T_p^{r_p})_{hG}$ has cells in degrees
$0,\ldots,b-1$ relative to $r_p(p-1)$. This assertion also applies
to a virtual representation, after shifting by its virtual
dimension. Since $E_b$ is $(b-2)$-connected, its Thom homology in
relative degrees less than $b-1$ agrees with group homology with
coefficient $\F_p(\det^{r_p})$.

For fixed $b$ and sufficiently large $p$, these group homology
groups consist only of degree zero when $r_p$ is even, and vanish
in this range when $r_p$ is odd. To see this, use the cyclic-group
generators $\eta,\xi$ from Lemma~\ref{layer:acyclic-sphere} and take
multiplier invariants; the corresponding homology calculation is
its degreewise dual. The sign of the multiplier permutation by
$\alpha$ is $\alpha^{(p-1)/2}$ modulo $p$, as follows by applying
the permutation to the Vandermonde product. Thus the odd
determinant twist shifts the required multiplier weight by
$(p-1)/2$. For a fixed degree range this weight, and every positive
untwisted invariant weight, eventually lies beyond that range.

In the even case, choose the identity vertex in the first join
factor at every stage. The fiber-point inclusions then commute with
the maps from stage $b$ to stage $b+4$, and hence induce maps of
their cofibers. Take these cofibers as remainders; in the odd case,
keep the entire partial Thom spectrum. For $b\geq2$ each remainder
is a finite $p$-local cell spectrum with relative cells in degrees
at most $b-1$, and it has no $p$-local homology below degree $b-1$.
Here mod-$p$ vanishing gives $p$-local vanishing because the homology
groups are finitely generated. Stable Hurewicz after desuspension
therefore gives the same lower connectivity bound. The remainder
at $b+4$ is connective in relative degree $b+3$, so the transition
from the remainder at $b$ to the remainder at $b+4$ is null by
comparison with the source cell dimensions. This null holds at
almost all primes for each fixed $b$.

The common virtual dimension shift does not affect the argument,
even if $r_p$ varies without bound. Moreover the transition remains
null after smashing with any $Z_p$. Since $Z_p$ has trivial action,
this smash commutes with the homotopy orbit construction. The
remainder telescope is consequently zero after the ultraproduct.
The compatible fiber-point maps in the even case give the asserted
equivalence.
\end{proof}

Write $\mathcal H(Z)=\map_{\cC}(\one,Z)$ for an ordinary sequence.
We next compute the whole $J^\vee$ coefficient object, with its
multiplication rather than only its homotopy groups.

\begin{lemma}[The coefficient algebra of the layer]
\label{layer:coefficient-algebra}
For every ordinary sequence $Z$, external multiplication is an
equivalence
\begin{equation}
 \label{layer:external-product}
 \map_{\cE}(\one,J^\vee)\otimes_D\mathcal H(Z)
 \xrightarrow{\simeq}
 \map_{\cE}(\one,J^\vee\otimes\Tr Z).
\end{equation}
The unit-input algebra is equivalent, as a commutative $D$-algebra,
to $D//\varpi$. Under the point and unit identifications of the
partial and full terms, the first map of
Equation~\eqref{eq:norm-triangle} for $J^\vee\otimes\Tr Z$ is
multiplication by $(p(p-1))_p$ on $\mathcal H(Z)$.
\end{lemma}

\begin{proof}
Dualizing the first Euler triangle gives
\[
 T^{-1}\longrightarrow\one\longrightarrow J^\vee.
\]
Lemma~\ref{layer:partial-thom} and exactness identify
$\operatorname{Part}(J^\vee\Tr Z)$ with $\mathcal H(Z)$, using the
fiber point and the unit of $J^\vee$. On the other hand,
\[
 (J_p^\vee\otimes\Tr Z_p)^{hG_p}
 \simeq\map\bigl(S(V)_{hG_p+},Z_p\bigr)
 \simeq Z_p
\]
by Lemma~\ref{layer:acyclic-sphere}. Evaluation at a point is the
inverse of the constant-class equivalence. By
Lemma~\ref{layer:norm-model}, norm followed by this evaluation on
the constant-unit copy sums all group actions. It is therefore
multiplication by $|G_p|=p(p-1)$.

To check Equation~\eqref{layer:external-product}, pair the norm
triangle for $J^\vee$ with $\mathcal H(Z)$ by tensoring trivial
maps. Do this first on the partial homotopy fixed point terms,
before using their norm identifications with homotopy orbits. These
pairings are balanced over $D$, acting by trivial sphere scalars,
and give a map of the entire join triangles. Passage to the
telescope gives the localization pairing on their cofibers.

The comparisons on the first two terms are equivalences. For the
partial term, the point-and-unit class at $Z=\one$, followed by
the norm, multiplies a class from $\mathcal H(Z)$ to the
corresponding point-and-unit-norm class for $Z$, by naturality at
each prime. Lemma~\ref{layer:partial-thom}, also applied to
suspended sphere classes, identifies this generator with the free
rank-one $D$-module and identifies its external product with
$\mathcal H(Z)$. For the full term the same assertion follows from
the constant generator. Exactness gives the equivalence on
cofibers.

At $Z=\one$, both initial terms are $D$ in degree zero, and their
map is multiplication by $\varpi(p-1)_p$. The factor $(p-1)_p$ is a
unit and $\varpi$ is a nonzerodivisor. The cofiber algebra thus has
homotopy only in degree zero, equal to $D/\varpi$, and its unit
induces the quotient map. Choose a null of $\varpi$ in this
algebra. The resulting algebra map from $D//\varpi$ is an
equivalence on underlying modules, hence an algebra equivalence.
\end{proof}

The coefficient computation identifies the proposed layer after
trivializing its source spheres. To compare it with the Euler
cofiber, we must still show that the restriction from $S(2V)$ to
$S(V)$ loses nothing. The following odd-twist vanishing supplies
that step.

\begin{lemma}[Odd twists vanish in the layer]
\label{layer:odd-vanishing}
For every ordinary sequence $Z$ and every sequence of odd integers
$r=(r_p)_p$,
\[
 \map_{\cE}(\one,T^rJ^\vee\Tr Z)=0.
\]
\end{lemma}

\begin{proof}
The dual Euler triangles give
$J^\vee\simeq\Sigma T^{-1}J$. It suffices to apply
Equation~\eqref{eq:norm-triangle} to $T^{r-1}J\Tr Z$ and show that
its first map is an equivalence.

At each prime, $S(V)$ has a finite equivariant cell decomposition
with prime-to-$p$ stabilizers. One may use the boundary of the
permutation simplex: a subgroup with $p$-torsion has no fixed point
in $S(V)$. Induction from each such stabilizer is $p$-locally a
retract of free induction by Lemma~\ref{power:prime-to-p}. The
ordinary full orbit-to-fixed-point norm on
$T^{r-1}J\Tr Z$ is therefore an equivalence, by its finite
equivariant cellular filtration.

Its full homotopy orbit spectrum is identified with the underlying
fiber by inclusion of a fiber point: $r-1$ is even, so use the
trivialization of Lemma~\ref{layer:even-trivializations}, followed
by the acyclic augmentation of
Lemma~\ref{layer:acyclic-sphere}. Its partial orbit spectrum is
identified by the same point. Indeed projection $J\to\one$ and
the Euler triangle reduce that assertion to the even calculation
for $T^{r-1}\Tr Z$ and the odd vanishing for $T^r\Tr Z$ in
Lemma~\ref{layer:partial-thom}. The projection from partial to
full orbits respects these fiber-point maps, so it is an
equivalence. Composing with the full norm proves the assertion.
\end{proof}

\subsection{The multiplicative comparison and completion}

\begin{proposition}[The multiplicative Euler layer]
\label{layer:comparison}
For every sequence of graded commutative algebras $X$, there are
natural product-compatible equivalences
\begin{equation}
 \label{eq:layer}
 B(X)//t\xrightarrow{\simeq}\Phi(X)
 \xleftarrow{\simeq}Q(X)\otimes_D(D//\varpi).
\end{equation}
The right-hand comparison applies to arbitrary spectrum inputs as
well. The scalar action in the $\delta$ direction agrees with
ordinary evaluation on $Q(X)$ followed by reduction. The
comparisons are compatible with algebra units and maps of inputs.
The first map uses the tautological Euler null on $S(2V)$, restricted
to $S(V)$.
\end{proposition}

\begin{proof}
For the right-hand equivalence, replace $W_i$ coherently by
$\Tr S(Fi)$ using Lemma~\ref{layer:even-trivializations} and
free-module adjunction. Apply
Equation~\eqref{layer:external-product} with the target shifted by
$S(-Fi)$. Multiplication of the coefficient algebra at index zero
with ordinary maps from these spheres gives the comparison in
every degree and weight. The same construction respects all
products, since the sphere trivializations are coherent and the
external pairing is multiplication in the mapping algebra. The
chosen agreement on the $\delta$ summand makes its scalar action
ordinary evaluation. Naturality in targets is built into all of
these maps. In particular no algebra structure on the target is
needed for the underlying coefficient comparison.

For the left-hand equivalence, first replace $J^\vee$ by
$J_2^\vee$ in the definition of $\Phi(X)$. In that algebra, null
$t$ by the tautological Euler null on $S(2V)$, multiplied by the
target unit. Concretely the tautological nonzero vector gives a
radial path from the zero section of the disk bundle to its
boundary, which becomes the point at infinity in the representation
sphere. This is the null in the dual Euler triangle
\[
 T^{-2}\longrightarrow\one\longrightarrow J_2^\vee.
\]
The universal property of the algebra quotient gives an algebra
map from $B(X)//t$ to this $J_2^\vee$ mapping algebra.

We verify its underlying map with the null specified, since the
existence of some module equivalence would not suffice. By
Lemma~\ref{found:free-null}, its source is the cofiber of
multiplication by $t$, and the map on the added module cell is the
chosen null multiplied by the input from $B(X)_{i-\tau}$. The
coherent identifications of $W_\tau$, $W_{i-\tau}$, and $W_i$
identify multiplication by $t$ with the map obtained from
$T^{-2}\to\one$ and the target unit. The multiplied null is,
by the same unitality, the geometric null in the dual Euler
triangle tensored with the target. The original input maps are
constant along $S(2V)$. Exactness consequently identifies this
particular map with the cofiber equivalence
\begin{equation}
 \label{layer:chosen-null-cofiber}
 \cofib\bigl(B(X)_{i-\tau}\xrightarrow{t}B(X)_i\bigr)
 \simeq\map_{\cE}
 \left(W_i,J_2^\vee\otimes
 \Tr\bigl(X_p(p\lambda+\delta)\bigr)_p\right).
\end{equation}

Now restrict along the inclusion of the first copy
$S(V)\to S(2V)$. Its quotient is stably $T\otimes J$:
the open complement with nonzero second component is an open
$V$-ball times $S(V)$, and its one-point compactification gives
this description. After dualizing, the fiber of restriction
therefore has an extra factor $T^{-1}J^\vee$. Every $W_i$ has
even twisting exponent in $V$, so moving it to the target leaves
an odd twist. Lemma~\ref{layer:odd-vanishing} makes this fiber
zero, including any ordinary sphere shift. Restriction is thus an
equivalence. Its composite with
Equation~\eqref{layer:chosen-null-cofiber} proves the first map
of Equation~\eqref{eq:layer}. All choices of nulls have been
multiplied with units and transported by the stated maps, so the
construction has the claimed unit and target compatibilities.
\end{proof}

\begin{corollary}[The first Koszul layer]
\label{layer:first-koszul}
For constant input $\one$, the comparison identifies
\[
 B(\one)//t\simeq F^*h_\Omega(Y_1)
\]
as graded algebras. For $M=(\MU_{(p)})_p$ there is a compatible
identification
\[
 B(M)//t\simeq F^*h_\Omega(M\otimes Y_1).
\]
These identifications preserve the algebra maps induced by the
unit and the scalar action in the $\delta$ direction.
\end{corollary}

\begin{proof}
For constant input $X$, $Q(X)=F^*h_\Omega(X)$. The first Koszul
attachment kills $b_0=\varpi u_0$. Its quotient is base change by
$D//\varpi$: send the null of $b_0$ to $u_0$ times the null of
$\varpi$. The chosen-null cofiber description and
Proposition~\ref{found:finite-cell-morita} give this identification also
after tensoring with $M$. Scalar base change supported in weight
zero is componentwise tensor on the other input, so it commutes
with the pullback of weights by $F$. Apply
Proposition~\ref{layer:comparison} and its unit compatibility.
\end{proof}

\begin{lemma}[Derived Euler completion]
\label{layer:completion}
For a rational graded commutative algebra $C$ with a homogeneous
chain-degree-zero polynomial class $t$, define
\[
 \widehat C=\lim_{s\geq1}C//t^s,
\]
using the natural derived polynomial quotient maps. For every
$r\geq1$, the natural map induces an equivalence
\[
 \widehat C//t^r\simeq C//t^r.
\]
Thus $\widehat C$ is complete for derived reduction by powers of
$t$. Any algebra extension whose underlying $\widehat C$-module
is perfect is complete as well. In particular these statements
apply to
$\widehat B(X)=\lim_s B(X)//t^s$.
\end{lemma}

\begin{proof}
Derived reduction modulo $t^r$ is tensoring over the polynomial
algebra with its two-term perfect quotient. It commutes with
limits. After this reduction, the cones of the maps from the
initial module $C$ to its quotients modulo $t^s$ form a pro-zero
tower: their transitions are multiplications by $t$ on the
corresponding shifts, and any $r$ consecutive transitions are
null because $t^r$ has its quotient nullhomotopy. Their inverse
limit is zero. Taking the limit of the reduced quotient maps
therefore gives the asserted equivalence. Applying this for all
$r$ proves completeness of $\widehat C$. A perfect
$\widehat C$-module is dualizable; tensoring it commutes with
this limit and with the finite quotient cofibers. The same
calculation proves the final assertion.
\end{proof}
